\scanpage{159}
% archive foliation: 158
Comme A est noethérien, il en est de même de $A'=A\{T_1,\ldots,T_r\}$ (7.5.4), donc $c)$ implique que B est noethérien; comme $\mathfrak J'=\mathfrak JA'$ est un voisinage ouvert de 0 dans $A'$ tel que les $\mathfrak J'^n$ forment un système fondamental de voisinages de 0, les images $\mathfrak J^nB$ des $\mathfrak J'^n$ forment un système fondamental de voisinages de 0 dans B, et comme on sait que B est séparé et complet, B est un anneau $\unsure{\mathrm J|\mathfrak J}B$-adique. Enfin, $B/\mathfrak JB$ est une algèbre (sur $A/\mathfrak J$) quotient de $A'/\mathfrak JA'=(A/\mathfrak J)[T_1,\ldots,T_r]$, donc est de type fini, ce qui achève de prouver que $c)$ entraîne $a)$.

Si B est $\mathfrak JB$-adique et noethérien, B est isomorphe à $\varprojlim B_n$, où $B_n=B/\mathfrak J^{n+1}B$ (7.2.11), et $\mathfrak JB/\mathfrak J^2B$ est un module de type fini sur $B/\mathfrak JB$. Soit $(a_j)_{1\leq j\leq s}$ un système de générateurs du $(B/\mathfrak JB)$-module $\mathfrak JB/\mathfrak J^2B$, et $(c_i)_{1\leq i\leq r}$ un système d'éléments de $B/\mathfrak J^2B$ dont les classes mod. $\mathfrak JB/\mathfrak J^2B$ forment un système de générateurs de la $(A/\mathfrak J)$-algèbre $B/\mathfrak JB$; on voit aussitôt que les $c_i a_j$ forment un système de générateurs de la $(A/\mathfrak J^2)$-algèbre $B/\mathfrak J^2B$, donc $a)$ implique $b)$.

Reste à prouver que $b)$ entraîne $c)$. L'hypothèse entraîne que $B_1$ est un anneau noethérien, et comme $B_1=B_2/\mathfrak J^2B_2$, on a $\mathfrak J^2B_1=0$, donc $\mathfrak JB_1=\mathfrak JB_1/\mathfrak J^2B_1$ est un $B_0$-module de type fini. Les conditions de (7.2.7) sont donc vérifiées par le système projectif $(B_n)$ et B est un anneau $\mathfrak JB$-adique. Soit $(c_i)_{1\leq i\leq r}$ un système fini d'éléments de B dont les classes mod. $\mathfrak J^2B$ engendrent la $(A/\mathfrak J^2)$-algèbre $B/\mathfrak J^2B$ (isomorphe à $B_1$).
\margin{\unsure{définition insuffisante} il faut aussi que les $c_i\in JB$ engendrent $JB/J^2B$ !}
Il existe donc un A-homomorphisme continu $u$ de $A'=A\{T_1,\ldots,T_r\}$ dans B qui se réduit à $\varphi$ dans A et est tel que $u(T_i)=c_i$ pour $1\leq i\leq r$ (7.5.3); si nous prouvons que $u$ est \textit{surjectif}, $c)$ sera établi, car de $u(A')=B$ on déduira $u(\mathfrak J^nA')=\mathfrak J^nB$, autrement dit $u$ sera un morphisme strict d'anneaux topologiques et B sera donc isomorphe à un quotient de $A'$ par un idéal fermé. Or, comme B est complet pour la topologie $\mathfrak JB$-adique, il suffit ([1], p. 18-07) de montrer que l'homomorphisme $\mathrm{grad}(A')\to\mathrm{grad}(B)$ déduit canoniquement de $u$ pour les filtrations $\mathfrak J$-adiques sur $A'$ et B, est surjectif. Mais par définition, les homomorphismes $A'/\mathfrak JA'\to B/\mathfrak JB$ et \uline{$\mathfrak JA'/\mathfrak J^2A'\to\mathfrak JB/\mathfrak J^2B$ déduits de $u$ sont surjectifs}; par récurrence sur $n$, on en déduit aussitôt qu'il en est de même de $\mathfrak JA'/\mathfrak J^nA'\to\mathfrak JB/\mathfrak J^nB$, et \textit{a fortiori} de $\mathfrak J^nA'/\mathfrak J^{n+1}A'\to\mathfrak J^nB/\mathfrak J^{n+1}B$, ce qui achève la démonstration.
\margin{!}

\textbf{7.6.} \textit{Anneaux complets de fractions.}

\textbf{(7.6.1)} Soient A un anneau linéairement topologisé, $(\mathfrak J_\lambda)$ un système fondamental de voisinages de 0 dans A formé d'idéaux, S une partie multiplicative de A. Soit $u_\lambda$ l'homomorphisme canonique $A\to A_\lambda=A/\mathfrak J_\lambda$, et pour $\mathfrak J_\mu\subset\mathfrak J_\lambda$, soit $u_{\lambda\mu}$ l'homomorphisme canonique $A_\mu\to A_\lambda$. Posons $S_\lambda=u_\lambda(S)$, de sorte que $u_{\lambda\mu}(S_\mu)=S_\lambda$. Les $u_{\lambda\mu}$ donnent canoniquement des homomorphismes surjectifs $S_\mu^{-1}A_\mu\to S_\lambda^{-1}A_\lambda$, pour lesquels ces anneaux forment un système projectif; désignons par $A\{S^{-1}\}$ la limite projective de ce système. Cette définition ne dépend pas du système fondamental de voisinages $(\mathfrak J_\lambda)$ choisi; en effet:

\textit{Proposition} \textbf{(7.6.2).} --- \textit{L'anneau $A\{S^{-1}\}$ est topologiquement isomorphe au séparé complété de l'anneau $S^{-1}A$ pour la topologie dont un système fondamental de voisinages de 0 est formé des $S^{-1}\mathfrak J_\lambda$.}

